\section*{ঐচ্ছিক Form ও Report অনুশীলন}

% Original MCQ 11

\item \begin{minipage}[t]{\linewidth}
ডেটা এন্ট্রির ইন্টারফেস কোনটি?
\par
\begin{tabularx}{\linewidth}{@{}X X@{}}
ক. ব্যাকআপ & খ. ফর্ম \\
গ. রিপোর্ট & ঘ. ইনডেক্স \\
\end{tabularx}
\end{minipage}

\par উত্তর: \textbf{খ।} ফর্ম ইনপুট ও সম্পাদনা সহজ করে।

% Original MCQ 12

\item \begin{minipage}[t]{\linewidth}
মুদ্রণযোগ্য ফলাফল তালিকার জন্য কোনটি উপযুক্ত?
\par
\begin{tabularx}{\linewidth}{@{}X X@{}}
ক. ডোমেইন & খ. রিপোর্ট \\
গ. ফরেন কী & ঘ. ইনডেক্স \\
\end{tabularx}
\end{minipage}

\par উত্তর: \textbf{খ।} রিপোর্ট তথ্য সাজিয়ে উপস্থাপন করে।

% Original MCQ 72

\item \begin{minipage}[t]{\linewidth}
ফর্মে বিভাগ তালিকা থেকে বাছাইয়ের কন্ট্রোল কোনটি?
\par
\begin{tabularx}{\linewidth}{@{}X X@{}}
ক. Label মাত্র & খ. Combo Box \\
গ. Report Header & ঘ. Page Footer \\
\end{tabularx}
\end{minipage}

\par উত্তর: \textbf{খ।} Combo Box বিকল্প মান বাছাইয়ের জন্য।

% Original MCQ 73

\item \begin{minipage}[t]{\linewidth}
এক শিক্ষার্থীর পরিচয়ের নিচে তার অনেক ফল দেখাতে কোনটি?
\par
\begin{tabularx}{\linewidth}{@{}X X@{}}
ক. Master Form ও Subform & খ. শুধু Page Header \\
গ. শুধু Ciphertext & ঘ. শুধু Label \\
\end{tabularx}
\end{minipage}

\par উত্তর: \textbf{ক।} 1:M তথ্যের উপযুক্ত ফর্ম বিন্যাস।

% Original MCQ 74

\item \begin{minipage}[t]{\linewidth}
রিপোর্টের প্রতি পৃষ্ঠায় কলামের শিরোনাম কোথায় রাখবে?
\par
\begin{tabularx}{\linewidth}{@{}X X@{}}
ক. Page Header & খ. Form Footer \\
গ. শুধু Detail & ঘ. Report Footer \\
\end{tabularx}
\end{minipage}

\par উত্তর: \textbf{ক।} Page Header প্রতি পৃষ্ঠার উপরে থাকে।

% Original MCQ 75

\item \begin{minipage}[t]{\linewidth}
সম্পূর্ণ রিপোর্টের সর্বমোট কোথায় উপযোগী?
\par
\begin{tabularx}{\linewidth}{@{}X X@{}}
ক. Page Header & খ. Report Footer \\
গ. Caption property & ঘ. প্রতি Detail-এ একই মান \\
\end{tabularx}
\end{minipage}

\par উত্তর: \textbf{খ।} রিপোর্ট শেষে সামগ্রিক সারাংশ রাখা যায়।

\textbf{৪. ফলের তালিকা ও ফর্ম।} একজন শিক্ষক শিক্ষার্থীর নম্বর Text Box-এ লিখছেন এবং বিভাগ Combo Box থেকে বেছে নিচ্ছেন। অধ্যক্ষ বিভাগভিত্তিক গড় নম্বর ও পুরো প্রতিষ্ঠানের সর্বমোট শিক্ষার্থীর সংখ্যা মুদ্রিত আকারে চান।
\begin{multicols}{2}
\begin{enumerate}
\item[ক.] \begin{minipage}[t]{\linewidth}Report কী?
\end{minipage}
\item[খ.] \begin{minipage}[t]{\linewidth}Label ও Text Box-এর কাজ আলাদা কেন?
\end{minipage}
\item[গ.] \begin{minipage}[t]{\linewidth}শিক্ষকের কাজের উপযুক্ত ফর্ম ও নম্বর যাচাইয়ের নিয়ম নির্ধারণ করো।
\end{minipage}
\item[ঘ.] \begin{minipage}[t]{\linewidth}অধ্যক্ষের প্রয়োজন মেটাতে কুয়েরি, গ্রুপিং ও রিপোর্ট সেকশনের ব্যবহার বিশ্লেষণ করো।
\end{minipage}
\end{enumerate}
\end{multicols}



\textbf{৪.} Label পরিচয় দেখায়, Text Box মান নেয়। Form-এ নম্বরের সীমা ও প্রয়োজনীয় ফিল্ডের নিয়ম বসাতে হবে। কুয়েরিতে বিভাগ অনুযায়ী AVG, রিপোর্টে Group Footer-এ বিভাগের সারাংশ এবং Report Footer-এ সামগ্রিক সংখ্যা রাখা যায়। শিরোনামের জন্য Page Header প্রয়োজন হতে পারে।



\subsubsection{ফর্ম: সহজ ডেটা এন্ট্রি}
\textbf{অনুশীলন:} (১) বিভাগ বাছাইয়ে Combo Box কেন সুবিধাজনক? (২) শিক্ষার্থীর পরিচয় ও বহু ফল দেখাতে কোন বিন্যাস ব্যবহার করবে? (৩) Label-এ ``মোবাইল'' লিখলে কি টেবিলে নতুন Mobile ফিল্ড তৈরি হয়?\par
\textbf{উত্তর-ইঙ্গিত:} নির্দিষ্ট তালিকা থেকে মান বাছাই; Master Form--Subform; হয় না।


\subsubsection{রিপোর্ট: তথ্যের উপস্থাপন}
\textbf{অনুশীলন:} (১) প্রতিটি পৃষ্ঠায় কলামের নাম কোথায় রাখবে? (২) প্রতিষ্ঠানের মোট বিক্রয় কোন সেকশনে উপযোগী? (৩) সব কর্মচারীর ঠিকানা একই মাপের কাগজে ছাপাতে কোন রিপোর্ট প্রয়োজন?\par
\textbf{উত্তর-ইঙ্গিত:} Page Header; Report Footer; Label Report।


